\documentclass[peerreview]{IEEEtran}
\usepackage{cite}
\usepackage{url}
\usepackage[utf8]{inputenc}
\usepackage{booktabs}
\usepackage{graphicx}
\usepackage{amsmath}
\usepackage{amsfonts}
\usepackage{amssymb}

\hyphenation{op-tical net-works semi-conduc-tor meta-fea-ture meta-fea-tures}

\begin{document}
\title{DataMetaMap: A Library for Dataset Vector Representation}

\author{Vladislav Minashkin, Ivan Papay, Vladislav Meshkov, Ilya Stepanov \\
\textbf{Intelligent Systems} \\
\textbf{MIPT 2025}}

\date{2025}

\maketitle
\tableofcontents
\listoffigures

\IEEEpeerreviewmaketitle

\begin{abstract}
This report presents \textit{DataMetaMap}, a Python library designed to 
represent multiple datasets in a unified vector space, enabling principled 
comparison of dataset similarity. The library offers a comprehensive suite 
of dataset embedding techniques compatible with PyTorch, addressing the 
challenge of transferring knowledge across datasets with varying schemas in 
meta-learning pipelines. We study four complementary approaches:
Maximum Mean Discrepancy (MMD), Task2Vec, Dataset2Vec, and
Wasserstein Task Embedding. The repository includes executable demos
and benchmark pipelines for dataset similarity analysis,
transfer-oriented retrieval.
\end{abstract}

\section{Introduction}
A fundamental challenge in modern machine learning is the efficient 
transfer of knowledge across tasks and datasets. Given a new target 
dataset, a practitioner must decide which pretrained model or which 
prior learning experience is most relevant — a decision that 
currently relies heavily on human intuition or exhaustive search.

A principled solution requires the ability to \emph{compare} datasets 
quantitatively: if datasets can be embedded into a common vector space 
where geometric proximity reflects task similarity, then finding the 
most relevant prior experience reduces to a nearest-neighbor query. 
Such embeddings, often called \emph{meta-features} or \emph{task 
embeddings}, form the foundation of data-driven meta-learning.

Existing approaches to dataset representation span a wide spectrum. 
Classical engineered meta-features (number of instances, class 
imbalance, statistical moments of features) are cheap to compute but 
have limited expressivity and require schema-specific design. 
Kernel-based methods such as Maximum Mean Discrepancy (MMD) provide
distribution-level comparisons without parametric assumptions but
scale poorly with dimensionality. More recent neural and
geometry-aware approaches --- Task2Vec, Dataset2Vec, and Wasserstein
Task Embedding --- learn rich, task-agnostic representations but
differ substantially in their inductive biases, computational cost,
and applicability 
to datasets with varying schemas.

We present \texttt{DataMetaMap}, a unified Python library that 
combines these dataset embedding methods within a practical,
PyTorch-compatible framework. Our contributions are:

\begin{itemize}
    \item A unified overview of four families of dataset comparison
    methods: MMD, Task2Vec, Dataset2Vec, and Wasserstein Task
    Embedding.
    \item A clean, extensible repository packaging Dataset2Vec and
    Wasserstein embedders together with a Task2Vec subpackage,
    shared utilities, and benchmark code.
    \item Experimental benchmarks for transfer-oriented dataset
    retrieval.
    \item Method-specific demo notebooks illustrating training,
    visualization, and similarity analysis in the learned embedding
    spaces.
\end{itemize}

\section{Notation and Conventions}
We adopt the following notation throughout this report.

\paragraph{Datasets and meta-datasets.}
A dataset is denoted $D = (X, Y)$, where
$X \in \mathbb{R}^{N \times M}$ is the predictor matrix and
$Y \in \mathbb{R}^{N \times T}$ is the target matrix; $N$, $M$, and
$T$ are the number of instances, predictors, and targets,
respectively. A \emph{meta-dataset} is a collection
$\mathcal{D}_{\mathrm{meta}} = \{D_1, \dots, D_K\}$ of (possibly
heterogeneous) datasets used for meta-training or evaluation.

\paragraph{Meta-features.}
A \emph{meta-feature extractor} is a map
$\varphi : \mathcal{D} \to \mathbb{R}^{d}$ that assigns to every
dataset a fixed-dimensional embedding $\varphi(D) \in \mathbb{R}^{d}$.
When the extractor is itself a learned (estimated) object, we write
$\hat{\varphi}$. For methods based on stochastic batch sampling,
$D_s$ and $D_s'$ denote multi-fidelity subsets (batches) drawn from
$D$ by subsampling instances, predictors, and targets.

\paragraph{Probe networks.}
For Task2Vec, we additionally introduce a \emph{probe network} with
parameters $\theta \in \mathbb{R}^{P}$ that is briefly fine-tuned on
the target dataset. The Fisher Information Matrix (FIM) of the probe
parameters is denoted $F(\theta) \in \mathbb{R}^{P \times P}$.

\paragraph{Distributional and kernel quantities.}
We write $k(\cdot, \cdot)$ for a positive-definite kernel (e.g.~RBF),
$\mu_P, \mu_Q$ for kernel mean embeddings of distributions $P$ and $Q$
in a reproducing kernel Hilbert space (RKHS) $\mathcal{H}$, and
$W_p(P, Q)$ for the $p$-Wasserstein distance. In the Wasserstein Task
Embedding section, $\psi$ denotes the MDS map used to embed
class-label distributions into a Euclidean space.

\paragraph{Hyperparameters.}
Throughout, $\gamma > 0$ denotes a bandwidth or scaling
hyperparameter for similarity models, and $\tau > 0$ a temperature
parameter where applicable. The binary indicator $i \in \{0, 1\}$
marks whether two batches originate from the same dataset ($i = 1$)
or not ($i = 0$).

\section{Proposed Solutions}

\subsection{Maximum Mean Discrepancy (MMD)}

\subsubsection*{Overview}
Gretton et al.~\cite{gretton2012kernel} introduced the Maximum Mean
Discrepancy as a non-parametric, kernel-based measure of the distance
between two probability distributions. Originally proposed as a
principled two-sample test, MMD operates by embedding distributions
into a reproducing kernel Hilbert space (RKHS) and measuring the
distance between their mean embeddings. The result is a scalar
dissimilarity score that can serve directly as a dataset distance
measure, without requiring an explicit parametric model.

\subsubsection*{Methodology}
Given two distributions $P$ and $Q$ with samples
$\{x_i\}_{i=1}^{m} \sim P$ and $\{y_j\}_{j=1}^{n} \sim Q$, the
squared MMD in an RKHS $\mathcal{H}$ induced by kernel $k$ is
\begin{equation}
    \mathrm{MMD}^2(P, Q)
    = \bigl\|\mu_P - \mu_Q\bigr\|_{\mathcal{H}}^{2},
\end{equation}
where $\mu_P = \mathbb{E}_{x \sim P}[k(x, \cdot)]$ is the kernel mean
embedding of $P$. An unbiased empirical estimator is
\begin{equation}
\begin{aligned}
    \widehat{\mathrm{MMD}}^2(P, Q)
    &= \frac{1}{m(m-1)} \sum_{i \neq i'} k(x_i, x_{i'}) \\
    &\quad + \frac{1}{n(n-1)} \sum_{j \neq j'} k(y_j, y_{j'}) \\
    &\quad - \frac{2}{mn} \sum_{i, j} k(x_i, y_j).
\end{aligned}
\end{equation}
For dataset comparison, $P$ and $Q$ are taken as the empirical
distributions of two datasets $D$ and $D'$, and
$\widehat{\mathrm{MMD}}^2(D, D')$ is used as their dissimilarity
score. A common choice is the RBF kernel
$k(x, y) = \exp\bigl(-\|x - y\|^{2} / (2\sigma^{2})\bigr)$ with
bandwidth $\sigma$ selected via the median heuristic.

\subsection{Task2Vec}

\subsubsection*{Overview}
Achille et al.~\cite{achille2019task2vec} propose Task2Vec, a method 
for embedding machine learning tasks into a fixed-dimensional vector 
space such that geometric distances between embeddings reflect 
meaningful task similarity. The core idea is to use the diagonal of 
the Fisher Information Matrix (FIM) of a pretrained \emph{probe
network} --- a fixed feature extractor briefly fine-tuned on the
target task --- as the task embedding. This captures the sensitivity
of the probe's parameters to the task's data distribution, producing 
embeddings that are stable, comparable across tasks, and correlated 
with transfer learning performance.

\subsubsection*{Methodology}
Given a dataset $D = (X, Y)$ and a probe network with parameters
$\theta$, the probe is fine-tuned on $D$ by minimizing the task loss
$\mathcal{L}(\theta; D)$. The Task2Vec embedding is the diagonal of
the empirical Fisher Information Matrix:
\begin{equation}
    \varphi(D)
    := \mathrm{diag}\bigl(\hat{F}(\theta)\bigr)
    \in \mathbb{R}^{P},
\end{equation}
where the empirical FIM diagonal is estimated as
\begin{equation}
    \hat{F}_{ii}(\theta)
    = \frac{1}{N} \sum_{n=1}^{N}
    \left(
        \frac{\partial \log p(y_n \mid x_n, \theta)}
             {\partial \theta_i}
    \right)^{2}.
\end{equation}
The resulting embedding $\varphi(D)$ encodes which parameters of the 
probe are most activated by dataset $D$. Task similarity is then 
measured via the cosine distance between embeddings:
\begin{equation}
    d(D, D')
    = 1 - \frac{\varphi(D)^{\top} \varphi(D')}
                {\|\varphi(D)\| \cdot \|\varphi(D')\|}.
\end{equation}

%------------------------------------------------------------
\subsection{Dataset2Vec}

\subsubsection*{Overview}
Jomaa et al.~\cite{jomaa2021dataset2vec} propose Dataset2Vec, a learned meta-feature extractor designed to replace hand-engineered dataset statistics in meta-learning pipelines. The key motivation is that classical engineered meta-features require substantial domain expertise and do not generalize well to datasets with different schemas, while autoencoder-based approaches are typically limited to fixed-schema data. Dataset2Vec addresses both issues by viewing a tabular dataset as a hierarchical set: a set of predictor and target variables, where each variable is itself a set of instance values. This representation is then processed by a permutation-invariant DeepSet architecture~\cite{zaheer2017deepsets}. To mitigate the lack of meta-labeled data, the authors introduce an auxiliary self-supervised task called \emph{dataset similarity learning}, which asks the model to determine whether two subsampled batches come from the same dataset.

\subsubsection*{Methodology}
Dataset2Vec represents a tabular dataset as a hierarchical set of predictor and target variables, where each variable is itself a set of instance values. Let
$D = (X^{(D)}, Y^{(D)}) \in \mathcal{D}$ denote a dataset, where
$X^{(D)} \in \mathbb{R}^{N(D)\times M(D)}$ and
$Y^{(D)} \in \mathbb{R}^{N(D)\times T(D)}$.
The meta-feature extractor is modeled as a DeepSet-based architecture
\begin{equation}
    \hat{\varphi} := h \circ g \circ f,
\end{equation}
where $f$ embeds individual predictor/target pairs, $g$ aggregates embeddings through pooling, and $h$ produces the final dataset representation. The resulting embedding can be written as
\begin{equation}
\hat{\varphi}(D)
:=
h\left(
\frac{1}{M(D)T(D)}
\sum_{m=1}^{M(D)} \sum_{t=1}^{T(D)}
g\left(
\frac{1}{N(D)}
\sum_{n=1}^{N(D)}
f\bigl(X^{(D)}_{n,m}, Y^{(D)}_{n,t}\bigr)
\right)
\right).
\end{equation}

To improve robustness and support datasets of varying sizes, embeddings are estimated from random multi-fidelity batches obtained by subsampling instances, predictors, and targets. The final dataset-level representation is computed as the average of multiple batch-level embeddings:
\begin{equation}
    \hat{\varphi}(D)
    := \frac{1}{B} \sum_{b=1}^{B}
    \hat{\varphi}\bigl(\mathrm{sample\text{-}batch}(D)\bigr).
\end{equation}

Training is performed using an auxiliary self-supervised task called dataset similarity learning. Given two batches $D_s$ and $D_s'$, the model predicts whether they originate from the same dataset using
\begin{equation}
    \hat{i}(D_s, D_s')
    := \exp\bigl(-\gamma\,\|\hat{\varphi}(D_s) - \hat{\varphi}(D_s')\|\bigr),
\end{equation}
and optimizes a symmetric binary cross-entropy objective over similar and dissimilar batch pairs.
\subsection{Wasserstein Task Embedding}

\subsubsection*{Overview}
Wasserstein Task Embedding (WTE) \emph{Liu et al. (2022)} is a
model-agnostic task representation method for supervised
classification that combines multidimensional scaling (MDS) with
Wasserstein embedding. The main idea is to represent every task as a
probability measure in an augmented space where the label information
is first embedded into a Euclidean vector space and then concatenated
with the input samples. This makes it possible to approximate a
hierarchical optimal transport dataset distance with a standard
Euclidean distance between task vectors.

Compared with pairwise optimal transport methods such as OTDD, WTE is
designed to be substantially faster when many tasks need to be
compared repeatedly. The paper shows that the resulting task
distances are strongly correlated with both forward transfer and
catastrophic forgetting, while also remaining highly correlated with
OTDD on benchmark task groups.

\subsubsection*{Methodology}
Let a task be given by a supervised dataset
$\tau = \{(x_n, y_n)\}_{n=1}^{N}$ with inputs $x_n \in \mathbb{R}^{d}$ and
labels $y_n \in \mathcal{Y}$. WTE first defines a label-to-label
distance matrix using the 2-Wasserstein distance between label
distributions. Following the simplification used in the paper, each
class-conditional label distribution $\nu_y$ is approximated by a
Gaussian, so that the distance between labels can be written in
closed form as the Bures--Wasserstein distance:
\begin{equation}
W_2^2(\nu_y, \nu_{y'})
=
\|u_y - u_{y'}\|_2^2
+
\operatorname{Tr}\!\left(
\Sigma_y + \Sigma_{y'}
- 2(\Sigma_y^{1/2}\Sigma_{y'}\Sigma_y^{1/2})^{1/2}
\right),
\end{equation}
where $u_y$ and $\Sigma_y$ denote the mean and covariance of the
Gaussian associated with label $y$.

The resulting pairwise label distances are then embedded into a
low-dimensional Euclidean space using multidimensional scaling.
Denoting the MDS map by $\psi$, each label $y$ is mapped to a vector
$\psi(\nu_y) \in \mathbb{R}^{l}$, where $l$ is chosen as a compromise
between accuracy and computational cost. This gives the
approximation
\begin{equation}
W_2^2(\nu_y, \nu_{y'})
\approx
\|\psi(\nu_y) - \psi(\nu_{y'})\|_2^2.
\end{equation}

After label embedding, each original sample-label pair is converted
to an augmented vector
\begin{equation}
[x, \psi(\nu_y)] \in \mathbb{R}^{d+l}.
\end{equation}
The original task is therefore transformed into an updated point
cloud in the augmented space, and the task distance becomes
\begin{equation}
d_{\tau}\bigl((x,y),(x',y')\bigr)
\approx
\left\|
[x, \psi(\nu_y)] - [x', \psi(\nu_{y'})]
\right\|_2.
\end{equation}
In this representation, the label discrepancy is absorbed into the
feature space, allowing the task comparison problem to be reduced to a
standard Wasserstein embedding problem.

To obtain the final task vector, WTE applies the Wasserstein
embedding framework with respect to a fixed reference measure
$\mu_0$. For each task distribution $\mu_i$ over the augmented
samples, an optimal transport map $T_i$ from $\mu_0$ to $\mu_i$ is
approximated, and the embedding is defined as
\begin{equation}
\Phi(\mu_i) = (T_i - \mathrm{id})\sqrt{p_0},
\end{equation}
or, in the discrete setting used in the implementation,
\begin{equation}
\Phi(X_i) = \frac{T_i - X_0}{\sqrt{N_0}},
\end{equation}
where $X_0$ denotes the reference sample set and $N_0$ is its size.
The Euclidean distance between the embedded task vectors then
approximates the 2-Wasserstein distance between tasks:
\begin{equation}
\|\Phi(\mu_i) - \Phi(\mu_j)\|_2 \approx W_2(\mu_i, \mu_j).
\end{equation}

Algorithmically, WTE proceeds in three steps: (1) compute the label
distance matrix, (2) embed labels with MDS, and (3) apply
Wasserstein embedding to the augmented task distributions. The paper
reports that this pipeline reduces the cost of task comparison from
pairwise optimal transport over all task pairs to a single transport
computation per task against the reference measure, which is the main
source of the computational gain.

%------------------------------------------------------------

\section{Research Methodology}
Our research methodology involved comprehensive implementation and 
experimental validation of each dataset embedding technique. We 
followed a systematic approach:

\begin{itemize}
    \item \textbf{Library Design}: We built a unified PyTorch-compatible 
    interface where each embedding method is implemented as a 
    self-contained module with consistent \texttt{fit()} and 
    \texttt{embed()} methods.
    \item \textbf{Algorithm Implementation}: Each method was implemented 
    as a separate class following a common base interface, enabling 
    straightforward benchmarking and extension.
    \item \textbf{Experimental Validation}: We conducted controlled 
    experiments on standard meta-dataset benchmarks to evaluate 
    embedding quality, dataset similarity retrieval accuracy, and 
    downstream meta-learning performance.
\end{itemize}

\subsection{Repository Structure}
The repository is organized into three main layers. The
\texttt{src/} directory contains reusable library components, the
\texttt{benchmarks/} directory contains executable evaluation
pipelines, and the \texttt{demo/} directory provides lightweight
notebooks for qualitative inspection of the methods. In addition, the
project includes unit tests for the main embedder modules and
continuous-integration workflows for automated testing and
documentation deployment.

\subsection{Benchmarking Pipelines in the Repository}
The public code base already contains concrete benchmark logic for
two practical use cases of dataset embeddings.

\paragraph{Task2Vec transfer-selection benchmark.}
The notebook \texttt{apply\_benchmark.ipynb}
uses Task2Vec embeddings to retrieve the nearest source dataset for a
given target image task. The benchmark operates on datasets such as
\texttt{mnist}, \texttt{cifar10}, \texttt{cifar100},
\texttt{letters}, and \texttt{kmnist}, computes embeddings with a
pretrained \texttt{ResNet-18} probe network, and then compares the
retrieval-based recommendation against logged pretrain-to-task
transfer results. The same notebook also defines random and large
pretraining baselines.

\paragraph{Wasserstein transfer-selection benchmark.}
The notebook
\texttt{benchmark\_wasserstein.ipynb}
implements a similar transfer-selection experiment using
Wasserstein-based dataset distances. It first computes class-wise
Bures--Wasserstein distances, derives dataset-level distances from the
resulting embeddings, and then recommends the closest source dataset
for each target task. When transfer logs are available, the notebook
compares these recommendations with observed downstream accuracies.

\paragraph{Status of MMD in the repository.}
MMD is retained in this report as an important classical baseline and
as part of the conceptual comparison framework. At the same time, the
current public repository places its strongest executable emphasis on
Dataset2Vec, Task2Vec, and Wasserstein components, including demos and
benchmark notebooks. Packaging MMD as a first-class module would be a
natural next step for future development.

\subsection{Demonstration Notebooks}
The repository also includes method-specific notebooks intended for
qualitative validation and onboarding.

\paragraph{Task2Vec demo.}
The notebook \texttt{demo/task2vec/simple\_example.ipynb} embeds a
small collection of image datasets with a pretrained
\texttt{ResNet-18} probe and visualizes the resulting task distances
with a clustered similarity matrix.

\paragraph{Dataset2Vec demo.}
The notebook \texttt{demo/dataset2vec/simple\_example.ipynb}
constructs a balanced synthetic tabular meta-dataset, trains a
Dataset2Vec encoder, and studies the resulting embeddings using
PCA/t-SNE projections together with nearest-centroid classification in
embedding space.

\paragraph{Wasserstein demo.}
The notebook \texttt{demo/wasserstein/simple\_example1 (1).ipynb}
demonstrates the full Wasserstein pipeline: computation of pairwise
class distances, MDS-based label embeddings, augmentation of the
feature space, and construction of reference-based task embeddings.

\section{Conclusions and Recommendations}
\texttt{DataMetaMap} provides a growing toolkit for dataset
representation and comparison. In its current public form, the
repository combines packaged Dataset2Vec and Wasserstein embedders, a
Task2Vec subpackage, theoretical coverage of MMD as a classical
baseline, and benchmark notebooks for transfer and meta-learning
scenarios.

The existing benchmark suite is designed to test the hypothesis that
learned embeddings capture richer dataset structure than purely
distributional baselines, especially when schemas vary across tasks.
We encourage further work on standardizing the benchmarking interface
across all methods, extending the released result tables with fully
reproducible numeric comparisons, and turning MMD into a packaged
executable component of the library.

\appendices
\section{Demo Experiments}
Our demo code is available at the GitHub repository. The experiments 
are divided into the following groups:

\begin{itemize}
    \item \textbf{Task2Vec Similarity Demo}: A notebook that embeds
    \texttt{mnist}, \texttt{cifar10}, \texttt{cifar100}, and
    \texttt{letters} with a pretrained \texttt{ResNet-18} probe and
    visualizes the resulting task-distance matrix.
    \item \textbf{Dataset2Vec Training and Visualization Demo}: A
    synthetic tabular meta-dataset is generated, a Dataset2Vec model
    is trained on balanced train/validation/test splits, and the
    learned representations are explored with PCA/t-SNE projections
    and nearest-centroid classification.
    \item \textbf{Wasserstein Embedding Demo}: A compact notebook
    computes class-wise Bures--Wasserstein distances, label embeddings,
    and final task embeddings for small image-dataset collections.
    \item \textbf{Transfer-Based Benchmark Notebooks}: Separate
    notebooks evaluate Task2Vec and Wasserstein retrieval as heuristics
    for choosing a source pretraining dataset before downstream
    fine-tuning.
\end{itemize}

\bibliographystyle{IEEEtran}
\bibliography{references}

\end{document}
